\documentclass[11pt]{article}

\usepackage[margin=1in]{geometry}
\usepackage[T1]{fontenc}
\usepackage[utf8]{inputenc}
\usepackage{hyperref}
\usepackage{longtable}
\usepackage{array}

\hypersetup{
  colorlinks=true,
  linkcolor=blue,
  urlcolor=blue
}

\newcommand{\code}[1]{\texttt{#1}}

\title{Datamak Multi-User Configuration}
\author{}
\date{Updated: 2026-03-26}

\begin{document}

\sloppy
\maketitle
\tableofcontents

\section{Purpose}

This document describes how Datamak was upgraded from a single-user repository
layout to a multi-user workflow. It covers both implementation details and the
operator workflow for configuring the application on a new account.

\section{Design Overview}

The main design goal is to remove assumptions that the active user is a
specific person or that runtime data always lives in a single personal
directory tree.

The upgrade introduces one shared configuration resolver in
\code{dtwin\_config.py} and one shell-oriented helper in
\code{tools/resolve\_dtwin\_env.py}. Active scripts now resolve settings in the
same order:

\begin{enumerate}
\item explicit CLI arguments,
\item environment variables,
\item user-local Datamak config,
\item derived defaults for standard Perlmutter and Flux layouts,
\item hard failure when a required value still cannot be inferred.
\end{enumerate}

\section{User-Local Config File}

\subsection{Location}

Datamak reads the user config from:

\begin{enumerate}
\item \code{DTWIN\_CONFIG}, if set,
\item \code{\$XDG\_CONFIG\_HOME/datamak/config.json}, if \code{XDG\_CONFIG\_HOME} is set,
\item \code{\$HOME/.config/datamak/config.json} otherwise.
\end{enumerate}

This keeps per-user workflow settings outside the repository checkout.

\subsection{Schema}

\begin{verbatim}
{
  "version": 1,
  "perlmutter": {
    "user": "",
    "host": "perlmutter.nersc.gov",
    "identity": "",
    "control_path": "/tmp/datamak_ssh_%r@%h_%p",
    "control_persist": "10m",
    "connect_timeout": 10,
    "monitor_timeout": 120,
    "base_dir": "",
    "batch_dir": "",
    "gx_path": ""
  },
  "flux": {
    "user": "",
    "host": "flux",
    "base_dir": "",
    "python_bin": ""
  }
}
\end{verbatim}

\section{Resolved Defaults}

\subsection{Perlmutter}

When Perlmutter settings are not explicitly configured, Datamak derives:

\begin{itemize}
\item user from \code{\$USER},
\item host as \code{perlmutter.nersc.gov},
\item base directory as \code{/pscratch/sd/<first-letter>/<user>/DTwin},
\item batch directory as \code{<base-dir>/newbatch},
\item GX path as \code{/global/homes/<first-letter>/<user>/GX/gx\_next6}.
\end{itemize}

\subsection{Flux}

When Flux settings are not explicitly configured, Datamak derives:

\begin{itemize}
\item user from \code{\$USER},
\item host as \code{flux},
\item base directory as \code{/u/<user>/DTwin/transp\_full\_auto},
\item Flux Python as \code{/u/<user>/pyrokinetics/.venv/bin/python}.
\end{itemize}

\subsection{Source metadata}

Source definitions such as Mate roots or TRANSP copy locations are not part of
the GUI-managed multi-user settings contract. They belong to the workflow data
model and should be represented through database records such as
\code{data\_origin}. The per-user settings window therefore does not edit them.

\section{Environment Variables}

The main environment variables are:

\begin{itemize}
\item \code{DTWIN\_CONFIG}
\item \code{DTWIN\_PERLMUTTER\_USER}, \code{DTWIN\_PERLMUTTER\_HOST},
\code{DTWIN\_PERLMUTTER\_BASE\_DIR}, \code{DTWIN\_PERLMUTTER\_BATCH\_DIR}
\item \code{DTWIN\_GX\_PATH}
\item \code{DTWIN\_FLUX\_USER}, \code{DTWIN\_FLUX\_HOST},
\code{DTWIN\_FLUX\_BASE\_DIR}, \code{DTWIN\_FLUX\_PYTHON}
\item existing SSH controls:
\code{DTWIN\_SSH\_IDENTITY}, \code{DTWIN\_SSH\_CONTROL\_PATH},
\code{DTWIN\_SSH\_CONTROL\_PERSIST}, \code{DTWIN\_SSH\_CONNECT\_TIMEOUT}
\end{itemize}

\section{GUI Workflow}

The Flask GUI still uses the route \code{/save\_hpc\_config} for compatibility,
but the saved data now lands in the user-local config file instead of
\code{db\_analysis/hpc\_config.json}.

The settings panel now captures:

\begin{itemize}
\item Perlmutter SSH and monitor settings,
\item Perlmutter remote base and batch directories,
\item GX installation path,
\item Flux SSH user, host, base directory, and Python path.
\end{itemize}

This means one person can configure the same repository checkout for their own
accounts without modifying repo-tracked files.

Source locations are intentionally excluded from this panel. Those belong to the
database workflow metadata rather than to user-local runtime settings.

\section{CLI and Shell Workflow}

\subsection{Python entrypoints}

Active Python entrypoints import \code{dtwin\_config.py} and resolve defaults
from the user-local config before falling back to derived layout conventions.

Examples include:

\begin{itemize}
\item Perlmutter batch deployment and transfer scripts in \code{batch/},
\item ingestion scripts such as
\code{db\_update/populate\_data\_equil\_from\_Mate\_KinEFIT.py},
\item the active Flux ingestion path under \code{db\_update/Transp\_full\_auto/}.
\end{itemize}

\subsection{Shell helper}

The shell helper can export resolved settings:

\begin{verbatim}
python3 tools/resolve_dtwin_env.py --profile perlmutter --format shell
python3 tools/resolve_dtwin_env.py --profile flux --format shell
\end{verbatim}

Typical usage in a shell script is:

\begin{verbatim}
eval "$(python3 "$DTWIN_ROOT/tools/resolve_dtwin_env.py" --profile flux --format shell)"
\end{verbatim}

\section{Flux Runtime Hand-Off}

The Flux full-auto wrapper now works in two stages:

\begin{enumerate}
\item the laptop-side launcher resolves the Flux profile and writes
\code{datamak\_runtime.env} into the upload bundle,
\item the Flux-side launcher and helper scripts source that file before running
Python entrypoints.
\end{enumerate}

This file contains exported Flux settings such as:

\begin{itemize}
\item \code{DTWIN\_FLUX\_USER}
\item \code{DTWIN\_FLUX\_HOST}
\item \code{DTWIN\_FLUX\_BASE\_DIR}
\item \code{DTWIN\_FLUX\_PYTHON}
\end{itemize}

As a result, the active Flux scripts no longer need hardcoded paths like
\code{/u/someuser/...}.

\section{Perlmutter Runtime Behavior}

Perlmutter batch deployment now resolves remote host, base directory, batch
directory, SSH controls, and GX path from the shared config layer.

The job submit scripts use:

\begin{verbatim}
GX_PATH="${DTWIN_GX_PATH:-$HOME/GX/gx_next6}"
\end{verbatim}

When a custom GX path is configured, the deploy scripts pass it into
\code{sbatch} through \code{--export=ALL,DTWIN\_GX\_PATH=...}.

\section{Migration from Legacy Repo-Local Config}

Older GUI settings were written to:

\begin{verbatim}
db_analysis/hpc_config.json
\end{verbatim}

The new config layer still reads that file as a legacy fallback if the
user-local config does not yet exist. Once a user saves settings through the
GUI, the new user-local config becomes the active source of truth.

This migration applies to Perlmutter and Flux runtime settings only. Source
definitions should remain database-backed rather than being migrated into the
GUI settings workflow.

\section{Typical Setup for a New User}

\begin{enumerate}
\item clone the repository,
\item create a Python environment and install the Datamak dependencies,
\item install or point to a working \code{pyrokinetics} environment,
\item open the GUI and save Perlmutter, Flux, and source-root settings,
\item verify SSH connectivity from the GUI using the test/open actions,
\item run the workflow scripts normally; they will resolve settings from the
user-local config.
\end{enumerate}

\section{Troubleshooting}

\begin{itemize}
\item If a script says a source root is missing, set the corresponding
\code{DTWIN\_*} variable or save it in the GUI settings panel.
\item If Perlmutter commands fail non-interactively, confirm the SSH identity
and control-path settings and ensure \code{sshproxy} is active if your site
requires it.
\item If Flux full-auto scripts cannot find Python, set
\code{flux.python\_bin} in the user config or export \code{DTWIN\_FLUX\_PYTHON}.
\item If an old checkout still appears to use repo-local settings, remove or
rename stale \code{db\_analysis/hpc\_config.json} after migrating to the
user-local config.
\end{itemize}

\section{Scope Note}

This upgrade covers the active GUI, Perlmutter batch flow, active ingestion
scripts, and \code{db\_update/Transp\_full\_auto/}. Older or duplicate trees
such as the top-level \code{transp\_full\_auto/} directory and
\code{db\_update/Transp\_full\_auto\_remote/} remain legacy audit targets and
should be cleaned up in a follow-up pass if they are reactivated.

\end{document}
